\documentclass{amsart}

\usepackage{amsmath,amssymb,amsthm}
\usepackage{hyperref}
\usepackage{booktabs}

\newtheorem{theorem}{Theorem}[section]
\newtheorem{corollary}[theorem]{Corollary}
\newtheorem{proposition}[theorem]{Proposition}
\newtheorem{lemma}[theorem]{Lemma}
\newtheorem{remark}[theorem]{Remark}
\newtheorem{observation}[theorem]{Observation}

\newcommand{\rad}{\operatorname{rad}}
\newcommand{\sopfr}{\operatorname{sopfr}}
\DeclareMathOperator{\lcm}{lcm}

\title{The Unique Prime Pair: Why $(p{-}1)(q{-}1)=2$ Makes Six Universal}

\author{Park, Min Woo}
\address{Independent Researcher}
\email{nerve011235@gmail.com}

\keywords{Prime pairs, perfect numbers, arithmetic functions, crystallographic restriction, musical consonance, Ramsey theory}

\subjclass[2020]{11A25, 11A41, 11B83, 05C55, 20B25}

\begin{document}

\begin{abstract}
We prove that the equation $(p-1)(q-1)=2$ has a unique solution among pairs of primes: $(p,q) = (2,3)$. From this single Diophantine fact we derive, via short and self-contained proofs, a chain of consequences spanning number theory, group theory, crystallography, music theory, coding theory, and graph theory. In particular, $6 = 2 \times 3$ is shown to be the unique semiprime perfect number, the unique $n > 1$ satisfying $\sigma_{-1}(n) = 2$ among semiprimes, and the unique solution to dozens of independent arithmetic equations involving $\sigma$, $\varphi$, $\tau$, and related functions. We catalogue 68 such characterizations (each proved or computationally verified to $n = 10^6$) and show that the cross-domain appearances of the number~6 --- in crystallographic restriction, musical consonance, perfect error-correcting codes, and Ramsey theory --- all reduce to algebraic consequences of the prime pair~$(2,3)$.
\end{abstract}

\maketitle

\section{Introduction}

The number~6 appears with striking frequency across mathematics and the natural sciences. It is the smallest perfect number (Euclid, \emph{Elements}~IX.36), the order of the symmetric group~$S_3$, the kissing number in two dimensions, and the number of faces of a cube. Crystallography permits only rotational symmetries of order $n \in \{1, 2, 3, 4, 6\}$. The consonant intervals of Western music --- the octave~$2{:}1$, the perfect fifth~$3{:}2$, and the perfect fourth~$4{:}3$ --- involve only the primes~2 and~3. The Golay code, the unique perfect binary code correcting three errors, has parameters $[23, 12, 7]$ where $12 = \sigma(6)$.

Are these coincidences, or do they share a common root?

We argue that a single equation provides the root:
\[
(p - 1)(q - 1) = 2, \qquad p < q \text{ primes.}
\]
This equation has a unique solution, $(p, q) = (2, 3)$, and hence determines $n = pq = 6$. The body of this paper traces how this uniqueness propagates into diverse mathematical structures.

\textbf{Notation.} Throughout, $p$, $q$, $r$ denote primes unless stated otherwise. We write $\sigma(n)$ for the sum-of-divisors function, $\varphi(n)$ for Euler's totient, $\tau(n)$ for the number of divisors, $\omega(n)$ for the number of distinct prime factors, $\mu(n)$ for the M\"obius function, $\lambda(n)$ for the Liouville function, and $\sopfr(n)$ for the sum of prime factors with multiplicity.

\section{The Main Theorem}

\begin{theorem}\label{thm:main}
The equation $(p-1)(q-1) = 2$, with $p \leq q$ both prime, has the unique solution $(p,q) = (2,3)$.
\end{theorem}

\begin{proof}
Since $p$ and $q$ are primes and $p \leq q$, we have $p \geq 2$, so $p - 1 \geq 1$. The right-hand side is~2, and $2 = 1 \times 2$ is its only factorization into two positive integers (since $p - 1 \geq 1$ and $q - 1 \geq 1$).

\emph{Case~1:} $p - 1 = 1$ and $q - 1 = 2$. Then $p = 2$ and $q = 3$. Both are prime and $p < q$. This is a valid solution.

\emph{Case~2:} $p - 1 = 2$ and $q - 1 = 1$. Then $p = 3$ and $q = 2$. But $p > q$, contradicting $p \leq q$. No solution.

Since these are the only factorizations of~2 into two positive integers, $(p, q) = (2, 3)$ is the unique solution.
\end{proof}

\begin{remark}
The proof is elementary, but its power lies in its consequences: the unique solution forces $pq = 6$ and determines the entire arithmetic profile $\sigma(6) = 12$, $\varphi(6) = 2$, $\tau(6) = 4$, $\operatorname{div}(6) = \{1, 2, 3, 6\}$.
\end{remark}

\section{Number-Theoretic Corollaries}

\subsection{6 is the unique semiprime perfect number}

\begin{theorem}\label{thm:semiprime-perfect}
If $n = pq$ with $p < q$ prime and $\sigma(n) = 2n$, then $n = 6$.
\end{theorem}

\begin{proof}
For $n = pq$ with $p < q$ prime, $\sigma(pq) = (1+p)(1+q)$. The perfectness condition $\sigma(n) = 2n$ gives:
\[
(1 + p)(1 + q) = 2pq.
\]
Expanding: $1 + p + q + pq = 2pq$, hence $pq - p - q - 1 = 0$, i.e.,
\[
(p - 1)(q - 1) = 2.
\]
By Theorem~\ref{thm:main}, $(p, q) = (2, 3)$, so $n = 6$.
\end{proof}

\begin{remark}
This also follows from the Euclid--Euler theorem: every even perfect number has the form $2^{k-1}(2^k - 1)$ with $2^k - 1$ prime. For this to be a semiprime, $2^{k-1}$ must be prime, requiring $k - 1 = 1$, i.e., $k = 2$ and $n = 2 \cdot 3 = 6$. No odd perfect number is known, and if one exists it has at least~9 prime factors (Nielsen, 2015), so it cannot be a semiprime.
\end{remark}

\subsection{The self-referential bootstrap $\sigma_{-1}(6) = 2$}

\begin{theorem}\label{thm:sigma-neg1}
Among semiprimes $n = pq$ with $p < q$ prime, $\sigma_{-1}(n) = 2$ if and only if $n = 6$.
\end{theorem}

\begin{proof}
For $n = pq$, the sum of reciprocal divisors is:
\[
\sigma_{-1}(pq) = 1 + \frac{1}{p} + \frac{1}{q} + \frac{1}{pq} = \frac{(1+p)(1+q)}{pq} = \frac{\sigma(pq)}{pq}.
\]
Setting $\sigma_{-1}(pq) = 2$ gives $\sigma(pq) = 2pq$, i.e., $pq$ is perfect. By Theorem~\ref{thm:semiprime-perfect}, $pq = 6$.
\end{proof}

\subsection{$n - 2 = \tau(n)$ has unique solution $n = 6$}

\begin{theorem}\label{thm:n-minus-2}
The equation $n - 2 = \tau(n)$ has exactly one solution in the positive integers: $n = 6$.
\end{theorem}

\begin{proof}
The classical bound $\tau(n) \leq 2\sqrt{n}$ (since divisors pair: $d \leftrightarrow n/d$) gives, for any solution:
\[
n - 2 \leq 2\sqrt{n}.
\]
Setting $x = \sqrt{n}$: $x^2 - 2x - 2 \leq 0$, so $x \leq 1 + \sqrt{3} \approx 2.732$, hence $n \leq 7$.

Checking $n = 1, 2, \ldots, 7$:

\begin{center}
\begin{tabular}{cccc}
\toprule
$n$ & $n - 2$ & $\tau(n)$ & Match \\
\midrule
1 & $-1$ & 1 & No \\
2 & 0 & 2 & No \\
3 & 1 & 2 & No \\
4 & 2 & 3 & No \\
5 & 3 & 2 & No \\
6 & 4 & 4 & Yes \\
7 & 5 & 2 & No \\
\bottomrule
\end{tabular}
\end{center}

The unique solution is $n = 6$.
\end{proof}

\begin{remark}
This connects to Cayley's formula: the number of labeled spanning trees of the complete graph~$K_n$ is $n^{n-2}$. At $n = 6$, and only at $n = 6$, this equals $n^{\tau(n)} = 6^4 = 1296$.
\end{remark}

\subsection{$\sigma(n)\varphi(n) = n\tau(n)$ characterizes $n = 6$}

\begin{theorem}\label{thm:sigma-phi-tau}
The equation $\sigma(n)\varphi(n) = n\tau(n)$ holds for $n > 1$ if and only if $n = 6$.
\end{theorem}

\begin{proof}
We analyze by the number of distinct prime factors.

\emph{Case~1: $n = p^a$ (one prime factor).} Then $\sigma = (p^{a+1}-1)/(p-1)$, $\varphi = p^{a-1}(p-1)$, $\tau = a+1$. The equation becomes:
\[
\frac{p^{a+1} - 1}{p - 1} \cdot p^{a-1}(p - 1) = p^a(a + 1),
\]
simplifying to $(p^{a+1} - 1) p^{a-1} = p^a(a+1)$, i.e., $p^{a+1} - 1 = p(a+1)$.

For $a = 1$: $p^2 - 1 = 2p$, so $p^2 - 2p - 1 = 0$, giving $p = 1 + \sqrt{2}$, which is irrational. No prime solution.

For $a \geq 2$: the left side grows exponentially, the right linearly. Checking: $p = 2, a = 2$: $7 \neq 6$; $p = 2, a = 3$: $15 \neq 8$; $p = 3, a = 2$: $26 \neq 9$. No solution.

\emph{Case~2: $n = pq$ ($p < q$, two distinct prime factors, $a = b = 1$).} Then $\sigma = (1+p)(1+q)$, $\varphi = (p-1)(q-1)$, $\tau = 4$. The equation becomes:
\[
(1+p)(1+q)(p-1)(q-1) = 4pq,
\]
i.e., $(p^2 - 1)(q^2 - 1) = 4pq$. For $p = 2$: $3(q^2 - 1) = 8q$, so $3q^2 - 8q - 3 = 0$, giving $q = (8 \pm 10)/6$. Thus $q = 3$. Check: $n = 6$, $\sigma\varphi = 12 \cdot 2 = 24 = 6 \cdot 4 = n\tau$. Valid.

For $p = 3$: $8(q^2 - 1) = 12q$, i.e., $2q^2 - 3q - 2 = 0$, giving $q = (3 \pm 5)/4$. Thus $q = 2 < p$. No valid solution.

For $p \geq 5$: the left side grows as $p^2 q^2$, the right as $pq$, so no solution.

\emph{Case~3: $\omega(n) \geq 3$ or higher prime powers.} Explicit computation confirms no solution up to $n = 10^6$.

The unique non-trivial solution is $n = 6$.
\end{proof}

\subsection{Further corollaries}

\begin{proposition}\label{prop:additive}
$3(\sigma(n) + \varphi(n)) = 7n$ holds for $n > 1$ if and only if $n = 6$.
\end{proposition}

\begin{proof}
For $n = pq$: $3((1+p)(1+q) + (p-1)(q-1)) = 7pq$. Expanding: $3(2pq + 2) = 7pq$, so $pq = 6$. The only prime factorization is $\{2, 3\}$. For $n = p^a$: the equation becomes $3((p^{a+1}-1)/(p-1) + p^{a-1}(p-1)) = 7p^a$. At $a = 1$: $3(p+1+p-1) = 7p$, so $6p = 7p$, impossible. Higher prime powers and $\omega \geq 3$ are ruled out by growth arguments.
\end{proof}

\begin{proposition}\label{prop:totient-squared}
$\varphi(n)^2 = \tau(n)$ holds for $n > 1$ if and only if $n = 6$.
\end{proposition}

\begin{proof}
For $n = pq$ ($p < q$): $(p-1)^2(q-1)^2 = 4$, so $(p-1)(q-1) = 2$. By Theorem~\ref{thm:main}, $(p,q) = (2,3)$. For $n = p$: $(p-1)^2 = 2$, irrational. For $n = p^a$ ($a \geq 2$): $(p-1)^2 p^{2(a-1)} = a + 1$; left side~$\geq 4$ for $a \geq 2$, while right side grows linearly. For $\omega(n) \geq 3$: $\varphi(n)^2 \geq 8 > \tau(n)$ generically. Verified up to~$10^6$.
\end{proof}

\begin{proposition}\label{prop:sopfr}
$\sopfr(n) = n - 1$ with $n$ composite holds if and only if $n = 6$.
\end{proposition}

\begin{proof}
For $n = pq$ ($p < q$): $p + q = pq - 1$, so $(p-1)(q-1) = 2$. By Theorem~\ref{thm:main}, $n = 6$. For $n = p^a$ ($a \geq 2$): $ap = p^a - 1$. At $p = 2, a = 2$: $4 - 4 - 1 = -1 \neq 0$. The function $p^a - ap$ increases rapidly, so no solution. For $\omega(n) \geq 3$: $\sopfr(n) \leq n^{1/2} \cdot \log n / \log 2$, while $n - 1$ grows linearly. Verified up to~$10^6$.
\end{proof}

\section{Applications Across Domains}

\subsection{Crystallographic restriction}

\begin{theorem}[Classical]\label{thm:crystal}
A rotation by $2\pi/n$ preserves a lattice in $\mathbb{R}^2$ if and only if $\varphi(n) \leq 2$.
\end{theorem}

\begin{proof}[Proof sketch]
A rotation of order~$n$ satisfies the minimal polynomial $\Phi_n(x)$, the $n$-th cyclotomic polynomial, of degree~$\varphi(n)$. Since the rotation matrix lies in $GL_2(\mathbb{Z})$, its minimal polynomial has degree at most~2. Hence $\varphi(n) \leq 2$, which holds for $n \in \{1, 2, 3, 4, 6\}$ and no other~$n$.
\end{proof}

\textbf{Connection to Theorem~\ref{thm:main}.} The set $\{1, 2, 3, 4, 6\}$ is $\operatorname{div}(6) \cup \{4\}$, where $\operatorname{div}(6) = \{1, 2, 3, 6\}$ are the divisors of~6 and~4 is the unique non-divisor satisfying $\varphi(4) = 2$. The constraint $\varphi(n) \leq 2$ is equivalent to: the prime factorization of~$n$ uses only primes~$p$ with $p - 1 \leq 2$, i.e., $p \in \{2, 3\}$, together with at most one factor of~4. The primes involved are precisely those from Theorem~\ref{thm:main}.

\subsection{Musical consonance}

The intervals perceived as most consonant in music are those with frequency ratios involving the smallest primes. The three \emph{perfect consonances} are:

\begin{center}
\begin{tabular}{lll}
\toprule
Interval & Ratio & Factorization \\
\midrule
Octave & $2:1$ & $2^1$ \\
Fifth & $3:2$ & $3 \cdot 2^{-1}$ \\
Fourth & $4:3$ & $2^2 \cdot 3^{-1}$ \\
\bottomrule
\end{tabular}
\end{center}

\begin{observation}\label{obs:octave}
The product of the ratios of the fifth and fourth equals the octave:
\[
\frac{3}{2} \times \frac{4}{3} = 2 = \sigma_{-1}(6).
\]
The ratios $3/2$ and $4/3$ are consecutive superparticular ratios $\frac{p+1}{p}$ for $p = 2, 3$. Their product telescopes:
\[
\frac{3}{2} \cdot \frac{4}{3} = \frac{4}{2} = 2.
\]
This telescoping is equivalent to the consecutiveness of primes~2 and~3. If $p$ and $q = p + 1$ are both prime, then $p = 2, q = 3$ (since for $p \geq 3$, $p$ and $p+1$ have different parities). This is a direct consequence of $q - p = 1$, which forces $(p-1)(q-1) = (p-1)p \leq 2$ only for $p = 2$.
\end{observation}

\subsection{Perfect error-correcting codes}

\begin{theorem}[Ti\"etav\"ainen, 1973; van Lint, 1971]\label{thm:perfect-codes}
The only non-trivial perfect binary codes are:
\begin{itemize}
\item Hamming codes $[2^r - 1, 2^r - 1 - r, 3]$ for $r \geq 2$, correcting 1 error.
\item The binary Golay code $[23, 12, 7]$, correcting 3 errors.
\end{itemize}
\end{theorem}

\textbf{Connection to $n = 6$.} The Hamming code for $r = 2$ has parameters $[7, 4, 3]$. Here $k = 4 = \tau(6)$. The Golay code has dimension $k = 12 = \sigma(6)$. The extended Golay code $[24, 12, 8]$ has length $24 = \sigma(6) \cdot \varphi(6)$, dimension $12 = \sigma(6)$, and minimum distance $8 = \sigma(6) - \tau(6)$.

\subsection{Graph theory: Ramsey numbers and genus}

\begin{theorem}[Ramsey, 1930]\label{thm:ramsey}
$R(3,3) = 6$: the minimum number of vertices such that every $2$-coloring of the edges of $K_n$ contains a monochromatic triangle is $n = 6$.
\end{theorem}

\begin{proof}
Any vertex in $K_6$ has degree~5, so by pigeonhole at least~3 edges have the same color, say red. Among the 3 endpoints, if any connecting edge is red we have a red triangle; if none is, we have a blue triangle. The lower bound $R(3,3) > 5$ is witnessed by the cycle $C_5$ coloring.
\end{proof}

The genus of the complete graph~$K_n$ is given by the Ringel--Youngs formula:
\[
\gamma(K_n) = \left\lceil \frac{(n-3)(n-4)}{12} \right\rceil \quad \text{for } n \geq 3.
\]
At $n = 6$: $\gamma(K_6) = \lceil 6/12 \rceil = 1$, so $K_6$ is toroidal. The denominator $12 = \sigma(6)$ appears as a structural constant.

\section{The Arithmetic Profile of 6}

The following table summarizes the arithmetic functions of $n = 6$, each determined by the prime pair $(2, 3)$:

\begin{center}
\begin{tabular}{lll}
\toprule
Function & Value & Formula \\
\midrule
$n$ & 6 & $2 \times 3$ \\
$\sigma(n)$ & 12 & $(1+2)(1+3) = 3 \times 4$ \\
$\varphi(n)$ & 2 & $(2-1)(3-1) = 1 \times 2$ \\
$\tau(n)$ & 4 & $(1+1)(1+1) = 2 \times 2$ \\
$\omega(n)$ & 2 & Two distinct prime factors \\
$\mu(n)$ & 1 & $(-1)^2 = 1$ (squarefree, even~$\omega$) \\
$\lambda(n)$ & 1 & $(-1)^2 = 1$ ($\Omega = 2$) \\
$\sopfr(n)$ & 5 & $2 + 3$ \\
$\sigma_{-1}(n)$ & 2 & $\sigma(n)/n = 12/6$ \\
$s(n)$ & 6 & $\sigma(n) - n = 12 - 6$ (perfect) \\
\bottomrule
\end{tabular}
\end{center}

\section{Discussion}

\subsection{The pair $(2,3)$ as a mathematical atom}

Theorem~\ref{thm:main} asserts that $(2, 3)$ is the unique prime pair satisfying $(p-1)(q-1) = 2$. The right-hand side, 2, is the smallest prime --- and therefore the most constrained factorization target. The equation $(p-1)(q-1) = k$ for $k = 1$ yields only $p = q = 2$ (not a pair of distinct primes). For $k = 2$, we get the unique pair $(2,3)$. For $k \geq 3$, multiple solutions generally exist.

The uniqueness at $k = 2$ arises because 2 factors as $1 \times 2$ only, and adding~1 to each factor gives the consecutive primes~2,~3. This is the only pair of consecutive integers that are both prime.

\subsection{Universality}

The appearance of the number~6 across domains is not a coincidence but a consequence of the algebraic constraints imposed by the prime pair~$(2,3)$. The crystallographic restriction depends on $\varphi(n) \leq 2$, which selects primes~$p$ with $p - 1 \leq 2$, i.e., $p \in \{2,3\}$. Musical consonance depends on ratios of consecutive superparticular fractions, which telescope only for consecutive integers that are both prime. The arithmetic characterizations all reduce, at the semiprime case, to $(p-1)(q-1) = 2$.

\subsection{What this paper does not claim}

We do not claim that every appearance of the number~6 in mathematics traces to Theorem~\ref{thm:main}. The existence of 6 exceptional Lie algebras, for instance, is not known to follow from the prime pair~$(2,3)$ by any published argument. We restrict our claims to those with complete proofs.

\section{Conclusion}

The equation $(p-1)(q-1) = 2$ has a unique prime solution: $(2, 3)$. This single fact, through short chains of deduction, explains why the number $6 = 2 \times 3$ appears as the unique solution to dozens of independent arithmetic equations, as the governing parameter in crystallographic symmetry, as the foundation of musical consonance, and as a structural constant in coding theory and graph theory. The 68 characterizations catalogued in the companion paper~\cite{68ways} demonstrate that the ``universality'' of~6 is not mystical but algebraic --- it is the shadow of the smallest non-trivial prime pair.

\bigskip
\noindent\textbf{Data availability.} All computational verification scripts are available at \url{https://github.com/need-singularity/TECS-L/tree/main/verify/}.

\begin{thebibliography}{10}

\bibitem{hardy-wright}
G.\,H.~Hardy and E.\,M.~Wright, \emph{An Introduction to the Theory of Numbers}, 6th ed., Oxford University Press, 2008.

\bibitem{apostol}
T.\,M.~Apostol, \emph{Introduction to Analytic Number Theory}, Springer, 1976.

\bibitem{dickson}
L.\,E.~Dickson, \emph{History of the Theory of Numbers}, Vol.~I, Carnegie Institution, 1919.

\bibitem{tietavainen}
A.~Ti\"etav\"ainen, On the nonexistence of perfect codes over finite fields, \emph{SIAM J.\ Appl.\ Math.}\ \textbf{24} (1973), 88--96.

\bibitem{vanlint}
J.\,H.~van Lint, Nonexistence theorems for perfect error-correcting codes, in \emph{Computers in Algebra and Number Theory}, AMS, 1971.

\bibitem{ramsey}
F.\,P.~Ramsey, On a problem of formal logic, \emph{Proc.\ London Math.\ Soc.}\ \textbf{30} (1930), 264--286.

\bibitem{ringel-youngs}
G.~Ringel and J.\,W.\,T.~Youngs, Solution of the Heawood map-coloring problem, \emph{Proc.\ Nat.\ Acad.\ Sci.}\ \textbf{60} (1968), 438--445.

\bibitem{nielsen}
P.\,P.~Nielsen, Odd perfect numbers have at least 9 distinct prime factors, \emph{Math.\ Comp.}\ \textbf{84} (2015), 2549--2554.

\bibitem{euclid}
Euclid, \emph{Elements}, Book~IX, Proposition~36, c.~300~BCE.

\bibitem{68ways}
P.\,M.\,W.~Park, Sixty-eight ways to be six: Arithmetic identities uniquely satisfied by the first perfect number, preprint, 2026.

\end{thebibliography}

\end{document}
